\documentclass[11pt]{article}
\usepackage[utf8]{inputenc}
\usepackage[T1]{fontenc}
\usepackage{lmodern}
\usepackage{geometry}
\geometry{a4paper, margin=1in}
\usepackage{hyperref}
\usepackage{xcolor}
\usepackage{titlesec}
\usepackage{parskip}
\usepackage{listings}

\definecolor{primary}{RGB}{0, 102, 204}
\titleformat{\section}{\Large\bfseries\color{primary}}{}{0em}{}[\titlerule]

\title{\vspace{-2cm}\textbf{\Huge 2048 AI Solving Competition}\\ \vspace{0.5cm} \Large Official Rules and Guidelines}
\author{}
\date{}

\begin{document}
\maketitle

\section{Introduction}
Welcome to the 2048 AI Solving Competition! Your objective is to build a Reinforcement Learning agent, search algorithm, or heuristic bot capable of mastering the classic game of 2048. Participants are provided with the complete, open-source Python environment to develop and train their models locally before submitting them for final evaluation.

\section{The Environment}
The official competition environment, \texttt{Game2048Env}, is given to you. The exact logic and random generation code used for evaluation are the same as the ones provided to you.
\begin{itemize}
    \item \textbf{State Space:} A 4x4 matrix representing the game board.
    \item \textbf{Action Space:} Discrete actions: \texttt{0 (Up), 1 (Right), 2 (Down), 3 (Left)}.
    \item \textbf{Rewards:} The reward for each step is the sum of the merged tiles. Invalid moves return a penalty of -1.
    \item \textbf{Dynamics:} Standard 2048 rules apply. After a valid move, a new tile spawns in a random empty cell (90\% chance of a 2, 10\% chance of a 4).
\end{itemize}

\section{Submission API and rules}
The submitted code has to be named \lstinline|agent_<yourname>.py| or \lstinline|agent_<yourname>.jl|. It needs the same main functions as provided to you \textcolor{red}{TBD HOW}.
Your submission needs some functions for us to test it easily:
\begin{itemize}
    \item \lstinline|train_and_instantiate()|: a function without arguments. When this function is called, the model is initialised and trained. returns a trained model.
    \item \lstinline|act(agent, state, valid_actions)|: a function that takes the trained model as an argument. When this function is called, it receives the current state of the game and the valid actions. It should return the action that the agent wants to take (a number between 0 and 3).
\end{itemize}

\section{Submission Limitations}
To ensure a standardized and automated evaluation process, all submissions must strictly adhere to the following limitations:
\begin{enumerate}
    \item \textbf{Language Constraints:} Agents must be written in \textbf{Python} (\texttt{.py}) or \textbf{Julia} (\texttt{.jl}). 
    \item Your agent may take at most 1 hour to train on a simple, CPU-only laptop. and \textcolor{red}{30 minutes, tbd} to execute. Training will be stopped at 1 hour, and the resulting agent will be used for scoring.
    \item \textbf{Zero-Argument Execution:} Your submitted file must be fully executable by calling the file directly \textbf{without any command-line arguments}.
    \begin{itemize}
        \item For Python submissions, the evaluation server will run: \texttt{python agent\_yourname.py}
        \item For Julia submissions, the evaluation server will run: \texttt{julia agent\_yourname.jl}
    \end{itemize}
    \item \textbf{Self-Contained Logic:} Your script must handle its own environment setup. It cannot load any pre-trained weights.
    \item \textbf{Dependencies:} You may use standard data science and machine learning libraries. For python these include PyTorch, TensorFlow, and everything in anaconda python version 3.19.3. In Julia, this includes Flux.jl EN NOG ANDERE TE BEPALEN . A list of pre-installed environment packages will be provided. If you have custom dependencies, include a \texttt{requirements.txt} or \texttt{Project.toml}.
\end{enumerate}

\section{Evaluation and Scoring}
Submissions will be evaluated on a secure server to prevent environment manipulation. 
\begin{itemize}
    \item \textbf{Trials:} Each AI will play 10 games, until the end.
    \item \textbf{Primary Metric:} Competitors will be ranked based purely on the \textbf{average final score} across all 10 games.
    \item \textbf{Tie-breaker:} In the event of identical average scores, the maximum tile achieved across the 100 episodes (e.g., reaching 8192 vs 4096) will determine the winner.
\end{itemize}

\section{TODO}
- 1 entry per person
- geen limiet om traintijd
- 100 games spelen in max 30 minuten
- model maximaal aantal MB
- indien speciale packages, overleg met Jury
- zelf trainen
- gemiddelde alle scores
- verwijs naar regels van het spel
- code van spel zelf geven (python en julia)
- code van test-environment geven (python en julia)
- maximum 2 files indienen (tenzij in overleg)
- elk om beurt 1 spel spelen. Zo live grafieken hebben om te updaten. (log as voor de punten)
- bijnamen verzinnen voor de modellen

Good luck, and may the best AI win!
\end{document}